% AST2000 Part 2: planetary orbits, Kepler's laws, and what an outside observer would see.
% Compile: cd "Project 2/paper" && latexmk -pdf part2.tex   (or: python build_papers.py 2 from the repo root)
% Every number written as \val{key} comes from results.tex, made by pro_2.py.
\documentclass[aps,prd,twocolumn,nofootinbib,floatfix,superscriptaddress]{revtex4-2}
\usepackage{../../common/paper/ast2000}

\begin{document}

\title{Part 2: Planetary Orbits, Kepler's Laws, and How Our Planets Look from Afar}
\author{Sakarias Frette}
\affiliation{AST2000 Rerun, University of Oslo}
\date{\today}

\begin{abstract}
We compute the orbits of the \val{n-planets} planets around our \val{star-mass}\,$\Msun$ star, both analytically and with a leapfrog integrator, over \val{n-home-orbits} orbits of the home planet.
Using the exact solution of Kepler's equation as the reference, we choose the time step from a measured error curve.
The worst position error is \val{worst-dev}\%, below the 1\% the official check allows, and \code{ast2000tools} accepted the orbits.
The simulation obeys Kepler's second law to $\val{area-diff-ecc}$ and the third law to $\val{kepler3-worst}$, and energy along the orbit satisfies the vis-viva law to $\val{visviva-err}$.
Letting the star move as well, we model what an observer far away would measure.
The heaviest planet, \val{c1-planet}, moves the star at \val{c1-star-speed}\,m/s, and from a noisy radial-velocity curve we recover $m\sin i$ to \val{c2-eye-err}\%.
Its transit, with a depth of \val{c3-depth}, gives the planet's radius and its density, \val{c3-rho}\,kg/m$^3$ (true: \val{c3-rho-true}).
A detectability survey shows that an observer could find the gas giants, but hardly any of the rocky planets.
As an extra, we measure general-relativistic perihelion precession: our code reproduces Mercury's \val{gr-mercury}$''$ per century.
\end{abstract}

\maketitle

% =====================================================================
\section{Introduction}
Before we can fly anywhere we need to know where the planets are, at any time, accurately enough to aim a spacecraft.
This part~\cite{ast2000project} builds that knowledge.
\begin{itemize}
\item Analytical and numerical orbits, verified by \code{ast2000tools}~\cite{ast2000tools} (\cref{sec:orbits}).
\item Kepler's laws as tests of the simulation (\cref{sec:kepler}).
\item The same system seen from outside: the star's wobble, radial velocities and transits (\cref{sec:wobble}).
\end{itemize}
We use astronomical units throughout: AU, yr and $\Msun$, in which $G=4\pi^2\AU^3\yr^{-2}\Msun^{-1}$.
Solar system seed: \val{seed}.

% =====================================================================
\section{Orbits}\label{sec:orbits}
\subsection{Theory}
For a planet that only feels a fixed star of mass $M$, the orbit is a conic section~\cite{ast2000notes},
\begin{equation}
  r(\theta)=\frac{a(1-e^2)}{1+e\cos(\theta-\omega)}, \label{eq:ellipse}
\end{equation}
with semi-major axis $a$, eccentricity $e$ and perihelion angle $\omega$.
The position at a given \emph{time} follows from Kepler's equation $M_{\rm anom}=E-e\sin E$.
Here the mean anomaly grows linearly in time, and $E$ is the eccentric anomaly.
We solve it with Newton's method for all planets and times at once.
It is our exact reference.

\subsection{Method}
We integrate all planets together with the kick--drift--kick leapfrog scheme,
\begin{align}
  \vb v_{i+1/2} &= \vb v_i + \tfrac12\vb a(\vb r_i)\Delta t,\nonumber\\
  \vb r_{i+1}   &= \vb r_i + \vb v_{i+1/2}\Delta t,\\
  \vb v_{i+1}   &= \vb v_{i+1/2} + \tfrac12\vb a(\vb r_{i+1})\Delta t.\nonumber
\end{align}
Leapfrog is symplectic~\cite{hairer2006}, so its energy error stays bounded instead of drifting.
For a central force it conserves angular momentum exactly.
Its error is a phase lag that grows linearly in time.
We measure it as the distance from the exact Kepler position, divided by the distance to the star, which is what the API checks.

\subsection{Results}
\Cref{fig:analytic,fig:overview} show the system.
The home planet orbits at $a=\val{home-a}\AU$ with $e=\val{home-e}$ and a period of \val{home-P}\,yr.
Periods range from \val{fastest-P}\,yr for the innermost planet to \val{slowest-P}\,yr for the outermost.
The three outer planets are gas giants with densities below 2\,g/cm$^3$.
The rocky planets have densities of 4--6\,g/cm$^3$.

The time-step study (\cref{fig:dt}) shows the expected $\Delta t^2$ convergence.
The minimum allowed 10\,000 steps per year would give a worst error of \val{dt-worst-10000}\%, which fails the check, because the innermost planet completes an orbit in only \val{fastest-P}\,yr.
We therefore use \val{steps-per-year} steps per year.
Over \val{sim-duration}\,yr this gives a worst error of \val{worst-dev}\% (planet \val{worst-dev-planet}, \cref{fig:poserr}).
Energy is conserved to $\val{energy-err}$ and angular momentum to $\val{h-err}$ (\cref{fig:conservation}).

\code{ast2000tools} accepted the orbits (\val{verified}).
Its own Kepler solver crashes at certain sample times, so we handed it every \val{api-thinning}. sample.
The API then saved the exact orbits, which all later parts use.

\widefig{a1_analytic_orbits}{Analytical orbits from \cref{eq:ellipse}, with the planets' positions at $t=0$.}{fig:analytic}
\widefig{a3_system_overview}{The system at a glance. Orbital speed against distance; mass against radius with lines of constant density; periods and eccentricities.}{fig:overview}
\widefig{a2_simulated_orbits}{Simulated orbits (color) on top of the analytical ellipses (dashed).}{fig:sim}
\colfig{a2_timestep_study}{Worst relative position error against the number of steps per year. Leapfrog converges as $\Delta t^2$.}{fig:dt}
\widefig{a2_position_error}{Position error against the exact Kepler solution for every planet: a phase lag that grows linearly.}{fig:poserr}
\widefig{a2_conservation}{Largest relative change in energy and angular momentum per time bin.}{fig:conservation}

% =====================================================================
\section{Kepler's laws}\label{sec:kepler}
\subsection{Second law}
The line from the star to a planet sweeps area at the rate $\dd A/\dd t=|\vb r\times\vb v|/2=h/2$.
This is constant because the angular momentum $h$ is conserved for any central force.
We measured the area swept in \val{area-fraction} of an orbit near perihelion and near aphelion (\cref{fig:kepler2}).
For the most eccentric planet (planet \val{ecc-planet}, $e=\val{ecc-e}$) it moves \val{speed-ratio-ecc} times faster at perihelion.
Yet the two areas agree to $\val{area-diff-ecc}$.

\subsection{Vis-viva}
Energy conservation along an ellipse gives $v^2=GM(2/r-1/a)$.
\Cref{fig:visviva} shows that the simulation follows it to $\val{visviva-err}$.
This is an independent check of the integrator.

\subsection{Third law}
Kepler's law with a fixed star, $P^2=a^3/M$, matches our measured periods to $\val{kepler3-worst}$ (\cref{fig:kepler3}).
Newton's version $P^2=a^3/(M+m)$ differs by a factor of about $1-m/2M$.
That is only \val{newton-corr}\% even for the heaviest planet (\val{newton-heaviest}), and our fixed-star model cannot see it.
The two-body simulation in \cref{sec:wobble} does.

\widefig{b1_kepler_second_law}{Kepler's second law. Equal areas swept in equal times near perihelion and aphelion (top). The areal velocity stays constant while the speed varies (bottom).}{fig:kepler2}
\widefig{b3_vis_viva}{Distance, speed and the vis-viva relation for the most eccentric planet.}{fig:visviva}
\widefig{b2_kepler_third_law}{Kepler's third law for all planets, and the size of Newton's correction.}{fig:kepler3}

% =====================================================================
\section{Seen from afar}\label{sec:wobble}
\subsection{The star's wobble}
The planet that moves the star the most has the largest $m/\sqrt a$.
Here that is planet \val{c1-planet}, with $m=\val{c1-mass}\Msun$ at $a=\val{c1-a}\AU$.
We integrate star and planet together with leapfrog in the centre-of-mass frame.
The star then moves on a small ellipse of radius \val{c1-wobble-km}\,km at \val{c1-star-speed}\,m/s (\cref{fig:wobble}).
The two-body energy $\tfrac12\hat\mu v^2-GM\hat\mu/r$ is conserved to $\val{c1-energy-err}$.
The measured period, \val{c1-P-sim}\,yr, matches Newton's $P^2=a^3/(M+m)$ (\val{c1-P-newton}\,yr).
Kepler's $a^3/M$ would give \val{c1-P-kepler}\,yr, which is \val{c1-kepler-days} days too long.

\subsection{Radial velocity}
An observer in the orbital plane's direction at inclination $i=\val{c2-incl}^\circ$ measures $v_r=v_{\rm pec}+v_{\star,x}\sin i$.
We add a peculiar velocity of \val{c2-vpec}\,m/s and Gaussian noise of \val{c2-noise}\,m/s, and take \val{c2-nobs} measurements.
We then analyse the curve as an outsider would.
The mean gives $v_{\rm pec}$, and a smoothed curve gives the period and amplitude.
For $m\ll M$ on a circular orbit, these give the minimum mass
\begin{equation}
  m\sin i = \frac{M^{2/3}\,v_\star\,P^{1/3}}{(2\pi G)^{1/3}}.
\end{equation}
By eye we read $P=\val{c2-P-eye}$\,yr and $v_\star=\val{c2-amp-eye}$\,m/s, so $m\sin i=\val{c2-msini-eye}\Msun$ (\cref{fig:rv}).
The true value is \val{c2-msini-true}, so we are off by \val{c2-eye-err}\%.
A least-squares sine fit gives \val{c2-msini-fit}.

\subsection{Transit}
Seen edge-on, the planet blocks a fraction $(R_p/R_\star)^2$ of the star's light.
We model the flux as a trapezoid with linear ingress and egress, add $\val{c3-noise}$ noise, and fit (\cref{fig:transit}).
The depth \val{c3-depth} gives $R_p=\val{c3-R-depth}$\,km.
Combined with the relative speed from the radial-velocity analysis, the \val{c3-ingress-h}-hour ingress gives \val{c3-R-ingress}\,km.
The true radius is \val{c3-R-true}\,km.
With the mass from radial velocity, the density is \val{c3-rho}\,kg/m$^3$ (true: \val{c3-rho-true}): a gas giant.

\subsection{Several planets}
With all three gas giants (planets \val{c4-planets}) in a full $N$-body run of \val{c4-duration}\,yr, the star's velocity is the sum of three wobbles (\cref{fig:multi}).
We search the noisy curve with a Lomb--Scargle periodogram~\cite{lomb1976}.
At each round we subtract the strongest sine and search again (``prewhitening''), and at the end we fit all three sines jointly.
The true periods are \val{c4-true-P}\,yr, and the joint fit gives \val{c4-fit-P}\,yr.
It recovers the two heaviest giants to within a few percent but misplaces the lightest one.
That planet (\val{c5-weakest-giant}) moves the star by only \val{c5-weakest-giant-K}\,m/s, much less than the noise per point.
Over a span of only a few orbits it also blends with the harmonics of the two big signals.
Separating it would need longer or less noisy observations.

\subsection{Could they find us?}
\Cref{fig:detect} puts every planet against the noise of our synthetic observations.
In radial velocity, only planets \val{c5-by-rv} stand out at $5\sigma$ after averaging.
The home planet moves the star by only \val{c5-home-K}\,cm/s.
In transit, planets \val{c5-by-transit} are deep enough once a hundred points in transit are averaged.
But an observer must also be lined up with the orbit, which for the home planet happens with probability $R_\star/a=\val{c5-home-prob}\%$.

\widefig{c1_star_orbit}{The star's wobble around the centre of mass, the two orbits, and energy conservation.}{fig:wobble}
\widefig{c2_radial_velocity}{Noisy radial-velocity curve, the smoothed curve an observer would read, and the extracted period and amplitude.}{fig:rv}
\widefig{c3_light_curve}{Light curve of planet \val{c1-planet} crossing the star, with the fitted trapezoid.}{fig:transit}
\widefig{c4_multi_planet_rv}{Radial velocity with three giants (top) and the prewhitening periodograms (bottom).}{fig:multi}
\widefig{c5_detectability}{Detectability of every planet by radial velocity (left) and transit (right), against the noise of our synthetic data.}{fig:detect}

% =====================================================================
\section{Extra: general relativity}
In Schwarzschild spacetime the orbit equation gains a term $3GMu^2/c^2$ (with $u=1/r$).
This is a central force $\vb a=-(GM/r^2)(1+3h^2/c^2r^2)\hat{\vb r}$, so leapfrog still conserves $h$ exactly.
GR rotates the orbit by about $10^{-7}$\,rad per orbit, which is smaller than leapfrog's own numerical precession.
We therefore run every orbit twice with identical steps, Newtonian and relativistic, and subtract, so the numerical error cancels.
We track the perihelion with the Laplace--Runge--Lenz vector.

The method reproduces Mercury's \val{gr-mercury}$''$ per century (\cref{fig:gr}).
For our innermost planet (\val{gr-fast-planet}) it gives \val{gr-fast-meas}$''$ per century (predicted: \val{gr-fast-pred}).
Over the 20 checked orbits, GR moves the home planet by \val{gr-home-km}\,km.
That is far less than the planet-planet forces a full $N$-body run adds (\cref{fig:gr2}).

That run also revealed something important.
Planets \val{nbody-pair} pass only \val{nbody-dist}\,AU apart at $t=\val{nbody-time}$\,yr, so the generated system is dynamically unstable.
The API, whose truth is Kepler ellipses, ignores this.

\widefig{r1_perihelion_precession}{Measured against predicted GR perihelion precession for all planets and Mercury.}{fig:gr}
\widefig{r2_size_of_corrections}{Size of the GR and $N$-body corrections compared with our numerical error.}{fig:gr2}
\colfig{r3_exaggerated_precession}{GR precession made visible by lowering $c$: the orbit becomes a rosette.}{fig:rosette}

% =====================================================================
\section{Discussion and conclusion}
Choosing the time step from a measured error, instead of a guess, was the key step.
The innermost planet sets the requirement, and it needs four times more steps than the minimum allowed.
Kepler's equation turned out to be valuable well beyond this part.
It gives exact positions at any time, and Parts 3--5 use it instead of interpolating our own orbits.

The outside view shows how biased planet detection is.
Heavy, close-in planets dominate radial velocity, large planets dominate transits, and Earth-like planets hide below the noise.
Our system would look like a three-giant system to most distant observers.

\begin{acknowledgments}
\code{pro_2.py} reproduces every figure and number; \code{rv_curve.txt} and \code{light_curve.txt} are the synthetic data for the partner group.
\end{acknowledgments}

\bibliography{../../common/paper/references}

\appendix
\section{Mission log}
\begin{missionbox}{Part 2 key numbers (seed \val{seed})}
Star mass & \val{star-mass}\,$\Msun$ \\
Home orbit $a$, $e$, $P$ & \val{home-a}\,AU, \val{home-e}, \val{home-P}\,yr \\
Steps per year & \val{steps-per-year} \\
Worst position error & \val{worst-dev}\% \\
API verification & \val{verified} \\
Kepler 2nd / 3rd law & \val{area-diff-ecc} / \val{kepler3-worst} \\
RV $m\sin i$ error & \val{c2-eye-err}\% \\
Transit density & \val{c3-rho}\,kg/m$^3$ \\
\end{missionbox}

\end{document}
